\documentclass[11pt]{article}
\usepackage[a4paper,margin=1in]{geometry}
\usepackage{amsmath,amssymb}
\title{Sharp Summary: Austin et al. (2021) -- D3PM}
\author{}
\date{}
\begin{document}
\maketitle

\section*{Problem and contribution}
The paper generalizes multinomial discrete diffusion by allowing \emph{structured} transition matrices, not only uniform corruption. The key claim is that forward-kernel design is a first-order modeling choice in discrete diffusion.

\section*{Forward process (exact formulas)}
For categorical states with one-hot row vectors $\mathbf x_t$:
\[
q(\mathbf x_t\mid \mathbf x_{t-1}) = \mathrm{Cat}(\mathbf x_t;\, \mathbf x_{t-1}\mathbf Q_t),
\qquad
\bar{\mathbf Q}_t = \mathbf Q_1\cdots \mathbf Q_t.
\]
Thus
\[
q(\mathbf x_t\mid \mathbf x_0)=\mathrm{Cat}(\mathbf x_t;\,\mathbf x_0\bar{\mathbf Q}_t).
\]
A closed-form posterior is available:
\[
q(\mathbf x_{t-1}\mid \mathbf x_t,\mathbf x_0)
=\mathrm{Cat}\!\left(
\mathbf x_{t-1};
\frac{\mathbf x_t\mathbf Q_t^{\top}\odot \mathbf x_0\bar{\mathbf Q}_{t-1}}
{\mathbf x_0\bar{\mathbf Q}_t\mathbf x_t^{\top}}
\right).
\]
This equation is the algebraic backbone of reverse modeling in D3PM.

\section*{Kernel design space}
The paper studies:
\begin{itemize}
\item uniform kernels,
\item absorbing/masking kernels,
\item discretized-Gaussian/local kernels,
\item embedding-distance kernels.
\end{itemize}
All are plugged into the same framework; stationary law is controlled by kernel structure.

\section*{Reverse parameterization and objective}
The reverse model is built via an $x_0$-prediction parameterization:
\[
p_\theta(\mathbf x_{t-1}\mid \mathbf x_t)
\propto
\sum_{\tilde{\mathbf x}_0}
q(\mathbf x_{t-1},\mathbf x_t\mid \tilde{\mathbf x}_0)
\tilde p_\theta(\tilde{\mathbf x}_0\mid \mathbf x_t).
\]
Training uses variational loss plus auxiliary denoising CE:
\[
L_\lambda = L_{\mathrm{vb}} + \lambda\,\mathbb E_{q(\mathbf x_0)q(\mathbf x_t\mid \mathbf x_0)}[-\log \tilde p_\theta(\mathbf x_0\mid \mathbf x_t)].
\]

\section*{Insight and limits}
Strong point: clean closed-form discrete algebra with practical kernel engineering.
Limitation: discrete-time design (many steps), no CTMC theory in this paper.

\section*{Relevance to this repo}
\textbf{Direct and central.}
It directly justifies the two notebook kernels (geometric/local and masking) and the step-3 training paradigm (learn reverse conditionals from noised-clean pairs).

\end{document}
